\documentclass[conference]{IEEEtran}
\IEEEoverridecommandlockouts
% The preceding line is only needed to identify funding in the first footnote. If that is unneeded, please comment it out.
%Template version as of 6/27/2024

\usepackage{cite}
\usepackage{amsmath,amssymb,amsfonts}
\usepackage{algorithmic}
\usepackage{graphicx}
\usepackage{textcomp}
\usepackage{xcolor}
\def\BibTeX{{\rm B\kern-.05em{\sc i\kern-.025em b}\kern-.08em
    T\kern-.1667em\lower.7ex\hbox{E}\kern-.125emX}}
\begin{document}

\title{Energy-Efficient LLM Inference via Adaptive Model Selection\\
}

\author{\IEEEauthorblockN{Yejean Kim}
\IEEEauthorblockA{\textit{University of Texas at Austin} \\
\textit{Department of Electrical and Computer Engineering}\\
yejean.kim@utexas.edu}
\and
\IEEEauthorblockN{Isha Jannu}
\IEEEauthorblockA{\textit{University of Texas at Austin} \\
\textit{Department of Computer Science}\\
ishajannu@utexas.edu}
}

\maketitle

\begin{abstract}
Large language models (LLMs) can provide strong performance across a
variety of tasks, but their inference workloads consume computational
resources and energy. Using a larger model for every request may be
unnecessary when a smaller or task-specialized model can provide
sufficient response quality. This work investigates the relationship
between model size, task specialization, inference energy consumption,
and response quality. Three instruction-tuned models are evaluated:
Qwen2.5-1.5B-Instruct, Qwen2.5-3B-Instruct, and
Qwen2.5-Coder-1.5B-Instruct. The models are compared across factual
question answering, summarization, and code generation. GPU energy is
estimated from NVIDIA Management Library (NVML) power measurements,
while inference latency, generated tokens, throughput, and
task-specific quality are recorded. Preliminary experiments were
conducted on an NVIDIA L4 GPU. The study uses these
measurements to examine whether model selection based on task
requirements can reduce inference energy without an unacceptable loss
in response quality. Further experiments on an NVIDIA GTX 1080 Ti are
planned to extend the evaluation to a second hardware platform.
\end{abstract}

\section{Introduction}
Large language models (LLMs) have become increasingly prevalent in
applications involving question answering, summarization, code
generation, and general-purpose assistance. As these models grow in
size and capability, the computational resources required to serve
inference requests have also increased. In particular, inference
places demands on GPU resources and consumes energy, raising concerns
about the efficiency and scalability of LLM-based applications.

A common approach to improving response quality is to use increasingly
large models. However, larger models do not necessarily provide
proportional improvements across all tasks. Some requests may be
handled adequately by smaller general-purpose models, while others
may benefit from larger models or models specialized for particular
tasks. This creates an opportunity to improve inference efficiency by
selecting models according to the requirements of individual requests.

This work investigates the relationship between model size, task
specialization, inference energy consumption, and response quality.
Three models are evaluated: Qwen2.5-1.5B-Instruct,
Qwen2.5-3B-Instruct, and Qwen2.5-Coder-1.5B-Instruct. These models
represent two general-purpose models of different sizes and a smaller
model specialized for code generation. Their performance is compared
across factual question answering, summarization, and code generation
tasks.

The primary objective is to determine whether selecting models based
on task requirements can reduce inference energy consumption without
substantially sacrificing response quality. To investigate this
objective, GPU energy consumption, inference latency, token
generation performance, and task-specific quality metrics are
measured. The resulting comparisons are used to explore the potential
benefits of adaptive model selection.

This study addresses the following research questions:

\begin{itemize}
    \item How do model size and task specialization affect
    LLM inference energy consumption and performance?
    \item What tradeoffs exist between inference energy
    consumption and response quality across different task categories?
    \item To what extent can adaptive model selection
    reduce energy consumption while maintaining acceptable response
    quality?
\end{itemize}

The contributions of this work are threefold. First, it provides an
empirical comparison of inference energy consumption and performance
across general-purpose and task-specialized LLMs. Second, it evaluates
the relationship between energy consumption and response quality
across multiple task categories. Third, it uses these comparisons to
investigate the feasibility of an adaptive model-selection strategy
that assigns requests to models based on their expected computational
cost and quality requirements.

\section{Background and Related Work}

\section{Methodology}

\subsection{Model Selection}

Three locally hosted instruction-tuned models are evaluated:

\begin{itemize}
    \item \textbf{Qwen2.5-1.5B-Instruct:} A general-purpose
    instruction-tuned model with approximately 1.5 billion parameters.
    \item \textbf{Qwen2.5-3B-Instruct:} A larger general-purpose
    instruction-tuned model with approximately 3 billion parameters.
    \item \textbf{Qwen2.5-Coder-1.5B-Instruct:} A 1.5-billion-parameter
    model specialized for coding-related tasks.
\end{itemize}

The model selection enables comparisons between models of different
sizes and between general-purpose and task-specialized models.
Comparing the two 1.5-billion-parameter models also provides an
opportunity to investigate whether task specialization can improve
performance without requiring a larger model.

\subsection{Benchmark and Task Categories}

The evaluation consists of three task categories:

\begin{enumerate}
    \item \textbf{Factual question answering:} Prompts requiring the
    model to identify or produce a correct factual answer.
    \item \textbf{Summarization:} Prompts requiring the model to
    generate a concise summary of provided information.
    \item \textbf{Code generation:} Prompts requiring the model to
    generate code that satisfies a specified programming problem.
\end{enumerate}

The preliminary benchmark contains 100 prompts per category, for a
total of 300 prompts. Each of the three models is evaluated on all
300 prompts, resulting in 900 inference runs.

\subsection{Energy Measurement}

GPU energy consumption is measured using NVIDIA's Management Library
(NVML), which provides access to GPU power telemetry. Rather than
relying on total system power measurements, this approach measures
GPU board power during inference.

GPU power is sampled at intervals of approximately 0.05 seconds while
each request is being processed. The measured power values are
integrated over the inference interval using the trapezoidal rule to
estimate energy consumption. For power samples $P_i$ recorded at times
$t_i$, the estimated energy consumption is

\begin{equation}
E = \sum_{i=1}^{n-1}
\frac{P_i + P_{i+1}}{2}(t_{i+1}-t_i),
\label{eq:energy}
\end{equation}

where $E$ is energy in joules, $P_i$ represents GPU power in watts,
and $t_i$ represents the corresponding measurement time in seconds.

This method provides an estimate of GPU energy attributable to each
inference request. It does not represent total system energy
consumption.

\subsection{Performance Metrics}

The evaluation considers three primary performance metrics:

\begin{itemize}
    \item \textbf{Inference latency:} The elapsed time required to
    generate a response, measured in seconds.
    \item \textbf{Generated tokens:} The number of tokens produced by
    the model for a request.
    \item \textbf{Token generation throughput:} The number of generated
    tokens per second, calculated by dividing the generated token
    count by inference latency.
\end{itemize}

Energy consumption and performance are analyzed together to evaluate
the computational cost of generating responses.

\subsection{Response Quality Evaluation}

Because the three task categories require different types of outputs,
task-specific evaluation metrics are used.

\begin{itemize}
    \item \textbf{Factual question answering:} Accuracy is measured by
    comparing the model's extracted multiple-choice answer against
    the reference answer.
    \item \textbf{Summarization:} ROUGE-L is used to measure overlap
    between generated summaries and reference summaries based on the
    longest common subsequence.
    \item \textbf{Code generation:} MBPP+ pass@1 is used to evaluate
    whether generated code passes the associated programming tests.
\end{itemize}

These metrics provide task-specific measures of response quality.
Their values are not directly interchangeable across task categories.

\section{Experimental Setup}

\subsection{Hardware and Software Environment}

The preliminary experiments were conducted using a Google Colab
environment equipped with an NVIDIA L4 GPU. The models
were loaded and executed locally within the environment using the
Hugging Face Transformers library and PyTorch.

GPU power measurements were collected through NVML. The experimental
implementation also used Python libraries for benchmark processing,
data storage, and response evaluation.

\subsection{Inference Configuration}

Each model was evaluated using the same benchmark prompts and a
consistent inference procedure. A system instruction was used to
encourage concise responses:

\begin{quote}
\texttt{Answer the user's question clearly and concisely.}
\end{quote}

The maximum input length was limited to 1,024 tokens. Maximum
generation lengths were configured by task category, as shown in
Table~\ref{tab:generation-limits}.

\begin{table}[t]
\caption{Task-Specific Maximum Generation Lengths}
\label{tab:generation-limits}
\centering
\begin{tabular}{lc}
\toprule
\textbf{Task Category} & \textbf{Maximum New Tokens} \\
\midrule
Factual question answering & 4 \\
Summarization & 128 \\
Code generation & 256 \\
\bottomrule
\end{tabular}
\end{table}

Generation was configured for deterministic decoding with sampling
disabled. The task-specific generation limits were used to constrain
inference workloads and reflect differences in expected response
length.

\subsection{Experimental Procedure}

For each benchmark prompt, the inference procedure was performed
independently for all three models. Before and after inference, GPU
synchronization was used to ensure that the measured execution
interval corresponded to completed model operations.

During each request, the experiment recorded inference latency,
generated token count, token generation throughput, GPU power samples,
and estimated GPU energy consumption. Generated responses were
subsequently evaluated using the corresponding task-specific quality
metric.

The resulting measurements were stored in CSV files for subsequent
statistical analysis. The preliminary dataset contains 900 inference
records, with 100 observations for each model--task combination.

\subsection{Preliminary and Final Evaluation}

The initial experiments were conducted on the Google Colab L4 GPU
to establish the measurement and evaluation procedure and obtain
preliminary comparisons among the selected models.

Additional experiments are planned on the UTCS Condor GPU cluster
using an NVIDIA GTX 1080 Ti. These experiments will provide a second
hardware configuration for evaluating the models and will serve as
the primary results for the final analysis.

The two experimental environments will be reported separately to
avoid attributing differences in hardware performance to model size
or task specialization. Comparisons across environments will be
interpreted with consideration for differences in GPU architecture
and hardware capabilities.

\section{Results}

\section{Discussion}

\section{Limitations and Future Work}

\section{Conclusion}

\section*{Acknowledgment}

\begin{thebibliography}{00}
\bibitem{b1} International Energy Agency, Energy and AI. Paris, France: IEA, Apr. 2025. [Online]. Available: https://www.iea.org/reports/energy-and-ai
\bibitem{b2} J. Fernandez, C. Na, V. Tiwari, Y. Bisk, S. Luccioni, and E. Strubell, “Energy considerations of large language model inference and efficiency optimizations,” in Proc. 63rd Annu. Meeting of the Association for Computational Linguistics (ACL), Vienna, Austria, 2025, pp. 32556–32569. doi: 10.18653/v1/2025.acl-long.1563.
\bibitem{b3} F. Caravaca, Á. Cuevas, and R. Cuevas, “From prompts to power: Measuring the energy footprint of LLM inference,” arXiv preprint arXiv:2511.05597, 2025. [Online]. Available: https://arxiv.org/abs/2511.05597
\bibitem{b4} A. Chatterji, T. Cunningham, D. J. Deming, Z. Hitzig, C. Ong, C. Y. Shan, and K. Wadman, “How people use ChatGPT,” NBER Working Paper No. 34255, Sep. 2025. doi: 10.3386/w34255.
\bibitem{b5} R. Bommasani et al., “Holistic evaluation of language models,” arXiv preprint arXiv:2211.09110, 2022. [Online]. Available: https://arxiv.org/abs/2211.09110
\bibitem{b6} Hugging Face, “Transformers documentation: Quickstart,” Hugging Face. [Online]. Available: https://huggingface.co/docs/transformers/quicktour. Accessed: Sep. 24, 2026.
\bibitem{b7} M. Chkir, S. M. Danish, J. Höll, and A. Asad, “Understanding the energy scaling of large language model inference across context lengths and attention architectures,” arXiv preprint arXiv:2608.25096, 2026. [Online]. Available: https://arxiv.org/abs/2608.25096
\bibitem{b8} Texas Advanced Computing Center, “Vista documentation,” University of Texas at Austin. [Online]. Available: https://docs.tacc.utexas.edu/hpc/vista/. Accessed: Sep. 24, 2026.
\bibitem{b9} S. Liu, ``Wi-Fi Energy Detection Testbed (12MTC),'' 2023, gitHub repository. [Online]. Available: https://github.com/liustone99/Wi-Fi-Energy-Detection-Testbed-12MTC
\end{thebibliography}

\end{document}
